\chapter{Tích hợp Trí tuệ Nhân tạo và An toàn Hệ thống}

\section{Mục tiêu và phạm vi AI}
Module AI trong dự án đóng vai trò là một trợ lý phân tích dữ liệu chuyên sâu (AI Analyst), hỗ trợ người dùng khám phá các mối liên hệ hoặc trích xuất số liệu từ tệp dữ liệu khí hậu nằm ngoài các tính năng tĩnh của dashboard. Về mặt công nghệ, hệ thống sử dụng mô hình ngôn ngữ lớn Google Gemini (\texttt{gemini-flash-latest}) truy cập qua API, kèm cơ chế dự phòng chạy ngoại tuyến (mock offline) khi không cấu hình khóa API nhằm bảo đảm quá trình trình diễn và chấm điểm cục bộ luôn hoạt động ngay cả khi mất kết nối. Tác vụ AI được giới hạn ở việc sinh mã truy vấn chỉ đọc -- dưới dạng câu lệnh SQL trên DuckDB hoặc đoạn mã Python/pandas tùy người dùng lựa chọn -- kèm giải thích ý nghĩa, và tuyệt đối không được phép can thiệp vào mã nguồn backend hoặc chỉnh sửa dữ liệu gốc.

\section{Nguyên tắc không thực thi ngầm}
Dự án tuân thủ nghiêm ngặt nguyên tắc minh bạch dữ liệu và duy trì sự kiểm soát của con người (Human-in-the-loop) đối với mọi hoạt động của hệ thống. Theo đó, mô hình trí tuệ nhân tạo không bao giờ được phép tự động hoặc âm thầm thực thi bất kỳ truy vấn SQL nào trực tiếp vào cơ sở dữ liệu. Mọi mã nguồn SQL do AI khởi tạo đều bắt buộc phải hiển thị một cách tường minh và trực quan trên giao diện người dùng ở trạng thái chờ duyệt. Hơn nữa, người dùng hoàn toàn có quyền can thiệp, chỉnh sửa cấu trúc hoặc nội dung cú pháp SQL này ngay trong trình biên soạn văn bản trước khi đưa ra quyết định thực thi. Hệ thống chỉ chính thức chạy truy vấn trên nền DuckDB sau khi nhận được lệnh xác nhận phê duyệt một cách chủ động từ phía người dùng qua hành động nhấn nút kích hoạt trên giao diện.

\section{Luồng phê duyệt của người dùng}
Quy trình thực thi truy vấn an toàn được thiết kế thông qua các bước phối hợp chặt chẽ giữa frontend và backend. Ban đầu, người dùng nhập câu hỏi bằng ngôn ngữ tự nhiên tại giao diện trợ lý AI, chẳng hạn như yêu cầu so sánh nền nhiệt độ cực đại giữa Sa Pa và Đà Lạt. Sau đó, giao diện gửi yêu cầu này tới API backend, truyền kèm ngữ cảnh cấu trúc bảng dữ liệu để dịch thành mã SQL thô. Khi nhận được kết quả dịch kèm giải thích logic từ AI Service, frontend sẽ hiển thị đề xuất này dưới dạng bản nháp tạm thời (DRAFT). Tại đây, người dùng tiến hành rà soát cấu trúc cú pháp câu lệnh, thực hiện các hiệu chỉnh cần thiết và quyết định nhấn nút duyệt chạy. Lúc này, hệ thống kích hoạt module kiểm soát an toàn để xác thực tính hợp lệ của câu lệnh, nếu vượt qua kiểm tra, DuckDB sẽ thực thi trực tiếp trên tệp Parquet cục bộ và trả kết quả dạng bảng hiển thị lên màn hình. Cuối cùng, toàn bộ quá trình tương tác sẽ được tự động ghi nhận vào cơ sở dữ liệu nhật ký kiểm toán.

\section{Thiết kế API}
Trách nhiệm giữa các API được phân định rõ ràng nhằm đảm bảo tính bảo mật và độc lập:
\begin{table}[htbp]
\centering
\renewcommand{\arraystretch}{1.25}
\begin{tabularx}{\textwidth}{|p{0.35\textwidth}|X|}
\hline
\textbf{API Endpoint} & \textbf{Trách nhiệm chi tiết} \\ \hline
\texttt{POST /api/ai/}\newline\texttt{proposals} & Tiếp nhận câu hỏi kèm ngữ cảnh lược đồ dữ liệu, gọi Gemini dịch thành mã nháp (SQL hoặc Python) và lời giải thích. Chưa thực thi. Trạng thái lưu: \texttt{draft}. \\ \hline
\texttt{POST /api/ai/}\newline\texttt{proposals/\{id\}/}\newline\texttt{approve} & Tiếp nhận mã SQL người dùng gửi lên sau khi duyệt/chỉnh sửa. Chạy bộ kiểm tra an toàn và lưu trạng thái \texttt{approved}. \\ \hline
\texttt{POST /api/ai/}\newline\texttt{proposals/\{id\}/}\newline\texttt{execute} & Chỉ cho phép chạy nếu đề xuất có trạng thái \texttt{approved}. Kết nối DuckDB thực thi SQL cục bộ, trả kết quả và lưu trạng thái \texttt{executed}. \\ \hline
\end{tabularx}
\caption{Phân tách trách nhiệm giữa các API AI}
\end{table}

\section{Hiển thị code và giải thích trên giao diện}
Giao diện người dùng được thiết kế để hiển thị đồng thời mã SQL đề xuất trong một khung biên soạn mã và một khung văn bản giải thích logic truy vấn:

\begin{figure}[htbp]
    \centering
    \includegraphics[width=0.82\textwidth]{dashboard_ai_analyst.png}
    \caption{Giao diện quản lý đề xuất và phê duyệt thực thi câu lệnh SQL của AI}
    \label{fig:dashboard-ai-analyst}
\end{figure}

\section{Kiểm soát an toàn truy vấn (Guardrails)}
Để bảo vệ an toàn cho cơ sở dữ liệu cục bộ, hệ thống được trang bị bộ lọc an toàn chuyên biệt để thực hiện xác thực nghiêm ngặt mã lệnh SQL trước khi chuyển giao cho động cơ truy vấn. Quy trình kiểm soát này được xây dựng trên ba lớp phòng ngự chính. Đầu tiên là cơ chế kiểm tra đặc quyền chỉ đọc, sử dụng thuật toán phân tích cú pháp để bắt buộc mọi câu lệnh gửi lên phải bắt đầu bằng các từ khóa truy xuất dữ liệu an toàn. Lớp phòng ngự thứ hai là danh sách cấm, ngăn chặn triệt để các câu lệnh có khả năng làm thay đổi tính toàn vẹn của dữ liệu hoặc phá hủy cấu trúc bảng. Cuối cùng, hệ thống thực hiện cơ chế chặn đa lệnh bằng cách chỉ cho phép thực thi duy nhất một truy vấn đơn lẻ trong mỗi phiên tương tác, loại trừ hoàn toàn nguy cơ chèn nhiều câu lệnh liên tiếp phân tách bằng ký tự đặc biệt. Đối với chế độ phân tích bằng Python/pandas, một bộ kiểm soát riêng phân tích cây cú pháp trừu tượng (AST) của đoạn mã để chặn triệt để mọi lệnh \texttt{import}, các hàm nguy hiểm như \texttt{eval}, \texttt{exec}, \texttt{open} cùng việc truy cập các thuộc tính đặc biệt (\enquote{dunder}); nhờ đó đoạn mã chỉ được phép tính toán thuần túy trên khung dữ liệu đã nạp sẵn mà không thể đọc/ghi tệp hay gọi mạng.

\section{Lưu trữ và khả năng kiểm toán (Audit Log)}
Toàn bộ nhật ký hoạt động của trợ lý trí tuệ nhân tạo được tự động ghi nhận vào cơ sở dữ liệu nhật ký kiểm toán chuyên dụng của hệ thống. Cơ chế ghi nhật ký được thiết kế để lưu giữ toàn diện các thông tin phục vụ mục đích giám sát hoạt động. Các thuộc tính lưu trữ bao gồm khóa định danh phiên tương tác duy nhất, nhãn thời gian thực thi, và nội dung câu hỏi ban đầu của người dùng. Ngoài ra, hệ thống cũng lưu trữ song song mã truy vấn do AI đề xuất ban đầu và mã truy vấn thực tế được người dùng phê duyệt chạy, cùng với trạng thái kết quả thực thi và thông tin chi tiết lỗi hệ thống trong trường hợp truy vấn gặp sự cố. Điều này đảm bảo tính minh bạch và cho phép quản trị viên dễ dàng rà soát hành vi của AI.

\section{Đánh giá chất lượng và Quá trình sử dụng}
Quá trình đánh giá thực nghiệm thông qua tập câu hỏi kiểm thử đa dạng cho thấy trợ lý AI hoạt động ổn định và chính xác. Các tác vụ phức tạp như so sánh lượng mưa giữa các khu vực địa lý, lọc nhiệt độ bất thường cục bộ hay tính toán hệ số tương quan đa biến đều được mô hình chuyển đổi thành mã SQL tối ưu cho DuckDB với thời gian xử lý cực ngắn. Hơn nữa, việc duy trì vai trò chủ động của con người giúp đảm bảo tính linh hoạt tối đa cho hệ thống. Trong trường hợp mô hình sinh ra câu lệnh có tên cột chưa khớp hoàn toàn hoặc bộ lọc chưa thực sự tối ưu, người dùng có thể chủ động sửa đổi trực tiếp trên giao diện để hoàn thiện câu lệnh trước khi thực thi, tạo nên cơ chế cộng tác hiệu quả giữa người dùng và trí tuệ nhân tạo.

